% ============================================================
%  THE ORTYX PROTOCOL
%  Part I — Immersion
%  Chapter 3: Harmonic Consensus and the Economy of
%             Structured Representation
% ============================================================

\chapter{Harmonic Consensus and the Economy of Structured Representation}
\label{ch:hcf1}

\chapepigraph{If the parts move together, you need only store the
fact of their moving together, not each part's motion separately.
The question is whether the world is structured enough for this
to be generally true.}{Flyxion}

\noindent
A harmonic series is not a collection of independent frequencies
that happen to share a common multiple. It is a relational structure:
the second harmonic is double the fundamental, the third is triple,
the fourth double the second. When a physical system produces a
harmonic series --- a vibrating string, a column of resonating air,
a vocal tract shaped to produce a vowel --- the amplitudes of the
individual harmonics evolve over time in ways that are not mutually
independent. They rise and fall together during the attack, decay
at rates governed by shared physical parameters, and modulate in
response to the same excitation source.

This relational structure is real. It is not a computational
convenience or an aesthetic preference. The harmonics of a vibrating
string are coupled because they share the same string. The harmonics
of a vocal sound are coupled because they are all driven by the
same glottal pulse train. The coupling is a physical fact, and it
means that a complete description of the amplitude evolution of each
harmonic independently contains enormous redundancy: most of what
is true about the third harmonic's envelope is already implied by
what is true about the fundamental's envelope.

The \HCFone{} architecture in the Phoenix corpus proposes to exploit
this redundancy systematically. Rather than encoding each harmonic
independently, the system stores the fundamental's amplitude envelope
as a master template and encodes each harmonic in terms of its
conformity to or deviation from that template. Harmonics that evolve
in lockstep with the fundamental require only a single bit to
describe: they conform. Harmonics that deviate require explicit
residual encoding. The expected result is a dramatic reduction in
the information required to represent a harmonic complex,
proportional to the degree of physical coupling between the harmonics.

\section{The Redundancy of Independent Harmonic Encoding}
\label{sec:harmonic-redundancy}

Standard spectral audio codecs encode audio by dividing the signal
into short frames, applying a transform (typically the modified
discrete cosine transform), and quantizing the resulting spectral
coefficients according to a psychoacoustic model. Within each frame,
the amplitude and phase of each spectral component are stored as
independent values. The independence assumption is computationally
convenient but discards structural information about how spectral
components relate to one another.

For harmonic signals --- which constitute a large fraction of
musical content --- this discarded structural information is
precisely what \HCFone{} proposes to encode explicitly instead.

Consider a sustained note on a cello at 220~Hz. The harmonics at
440, 660, 880~Hz and beyond do not evolve independently: the
bow-string interaction drives all of them simultaneously. If one
knows how the fundamental's amplitude is evolving, one can predict
the evolution of most other harmonics with considerable accuracy.
Standard codecs spend bits encoding this predictable variation.

\begin{technote}[Technical Note: Harmonic Conformity Relation]
Let the fundamental frequency be $f_0$ with amplitude envelope
$A_0(t)$, and let the $n$-th harmonic have amplitude envelope
$A_n(t)$. Define the normalized amplitude velocity:
\[
  v_n(t) = \frac{dA_n/dt}{A_n(t)}.
\]
The harmonic conformity relation is:
\[
  \Gamma_n(t) = \mathbf{1}\!\left(|v_n(t) - v_0(t)| < \varepsilon\right),
\]
where $\varepsilon$ is a conformity threshold. The sparse consensus
encoding is then:
\[
  \mathcal{H}_{\mathrm{cons}}
  = \bigl\{A_0(t),\, f_0,\, \phi_0,\,
    \Gamma_1, \ldots, \Gamma_N\bigr\},
\]
with residual deviations stored only when $\Gamma_n(t) = 0$.
If $k$ harmonics are non-conforming, the encoding cost is
$O(1) + O(k)$, compared to $O(N)$ for independent storage.
\end{technote}

The efficiency gain scales with the physical coupling of the
instrument. For highly coupled sources --- bowed strings, wind
instruments in steady state, sustained vocal vowels --- the
fraction of non-conforming harmonics is expected to be small
and the compression advantage large. For less coupled sources
--- percussion, plucked string attacks, complex mixtures ---
the conformity rate is lower and the advantage diminishes.

\section{Connection to Predictive Coding and Sparse Residual Encoding}
\label{sec:predictive-coding}

The \HCFone{} framework did not invent residual encoding. Predictive
coding, dating to the 1950s in the communications engineering
literature, encodes signals as the difference between a predicted
value and the actual value. Linear predictive coding of speech,
differential pulse-code modulation, and motion compensation in
video codecs all exploit this principle.

The \HCFone{} contribution is the specific application of the
conformity relation to the harmonic structure of musical audio ---
a domain-specific optimization that takes advantage of physical
facts about musical instruments that general-purpose predictive
coding does not explicitly model.

\begin{epistemicstatus}{Operationally Grounded}
The efficiency advantage of harmonic consensus encoding over
independent harmonic encoding is derivable from the conformity
relation and scales with the physical coupling of the source.
For strongly coupled harmonic sources, the theoretical advantage
is significant and testable with existing datasets and codec
implementations. The claim is well-grounded at \LevelI{}
(engineering utility) and connects to established literature
at \LevelII{} (computational metaphor, via predictive coding).
\end{epistemicstatus}

\section{The Buddy System and Envelope Sharing}
\label{sec:buddy-system}

The corpus extends the conformity relation into a broader
architectural principle: a hierarchical structure in which
harmonics that consistently track one another are grouped into
families, with the most fundamental component storing the shared
envelope and remaining members storing only their conformity
status and phase offsets.

The buddy system has intuitive appeal beyond its compression
application. It provides a natural account of how a listener
might group harmonic components into a unified percept: the
harmonics that ``move together'' are the harmonics that ``belong
together.'' Bregman's auditory scene analysis framework has
established that harmonic coherence is one of the primary cues
the auditory system uses to group spectral components into streams.
The \HCFone{} conformity relation operationalizes ``evolving
together'' in a specific, measurable way.

\begin{dangerzone}[Representation Drift]
The move from ``harmonics that are efficiently encoded together''
to ``harmonics that the auditory system groups together'' is a
transition from an engineering claim to a perceptual one. These
claims may align, and the alignment would be scientifically
interesting if confirmed. But encoding efficiency is a property
of the representation; perceptual grouping is a property of the
listener. Demonstrating their agreement for studied cases does
not establish that the conformity relation tracks the auditory
grouping mechanism. The distance between them is exactly where
representation drift begins.
\end{dangerzone}

\section{Compression as Discovery of Invariant Structure}
\label{sec:compression-invariant}

The \HCFone{} materials propose a reframing of what compression
is. Conventional lossy compression is framed as trade-off:
information is discarded in exchange for reduced cost.
The frame is inherently about loss.

\HCFone{} proposes instead: compression is the discovery of
invariant relational structure. A harmonic complex representable
as a fundamental plus conformity flags has not had information
removed; its genuine information content has been identified and
represented directly. The redundancy eliminated was never real
information --- it was the artifact of treating structurally
dependent quantities as independent.

\claimlabeltext{Level~I}{The compression efficiency advantage of
consensus encoding is real, derivable, and proportional to the
conformity rate. This is a solid engineering result.}

\claimlabeltext{Level~II}{The reframing of compression as
``discovery of invariant relational structure'' rather than
``discarding of information'' is a productive metaphor connecting
to minimum description length frameworks and Kolmogorov complexity.}

\claimlabeltext{Level~III}{The phenomenological claim that
listeners experience compressed audio as impoverished --- that
something genuinely heard is absent rather than measurably
different --- connects to this framing but is not derivable
from the information-theoretic account alone.}

\section{The Universality Temptation}
\label{sec:hcf1-universality}

The same pattern observed in Chapters~1 and~2 appears here.
Having established harmonic consensus at the level of audio,
the corpus begins extending it. The conformity relation, originally
defined for harmonic amplitude envelopes, is proposed as a general
model of how any system whose components are physically coupled
can be encoded efficiently.

The cautious version --- that the buddy system might model how
neural populations encode correlated activity --- is scientifically
tractable. The bolder version --- that ``cognition is in the
business of discovering invariant structure'' and consensus encoding
is its universal implementation --- is a claim at \LevelIII{}
or \LevelIV{} that the audio compression work does not establish.

It requires an account of how cognitive phenomena that appear to
involve active generation of structure --- imagination, planning,
counterfactual reasoning, creative synthesis --- fit within the
consensus encoding framework, or an argument that they do not
pose a genuine challenge. The corpus does not provide this.

The pattern appearing for the third consecutive time is worth
marking explicitly: a framework grounded at \LevelI{} makes
a genuine engineering contribution, demonstrates real explanatory
contact at \LevelII{}, and then begins reaching toward
\LevelIII{} and \LevelIV{} claims its foundation does not
directly support. This is not a failure mode yet. It is a research
programme generating ambitions that exceed its current results.
The question is whether those ambitions are structured to be
testable.

\section{The Three Frameworks as a System}
\label{sec:three-frameworks}

The \Marine{} is the perceptual gateway: temporal coherence
measured as jitter. \MEMeight{} is the cognitive layer: salience-filtered
events encoded as dynamic interference. \HCFone{} is the
representational layer: efficient encoding of structured signals
feeding both the gateway and the memory system.

Together, they constitute a \emph{salience-centered computational
paradigm}: an alternative to both symbolic AI and deep learning
grounded in temporal coherence, interference dynamics, and
relational compression.

\begin{epistemicstatus}{Reconstructive Conjecture}
The unified paradigm description is a reconstruction --- the
strongest reading of what the corpus seems to be building toward.
The individual components are developed separately in the source
materials; the integration is more programmatic than implemented.
\end{epistemicstatus}

The shared primitive --- temporal coherence --- runs through all
three layers. The \Marine{} detects it in incoming signals;
\MEMeight{} preserves it in memory through resonance; \HCFone{}
exploits it in compression through the conformity relation.
Whether this unity is productive or constraining depends on whether
temporal coherence is as central to cognition as the framework
assumes.

That tension is the conceptual fault line on which Part~II
will apply pressure. The framework is internally coherent,
its shared primitive is real, and its unification ambition
is genuine. The question is whether the real world is structured
enough around temporal stability to vindicate the ambition.
That question cannot be answered from inside the framework.

\medskip

Chapter~4 examines the AI architecture implications of the
salience-centered paradigm: the proposal that the
Marine--\MEMeight{}--\HCFone{} stack constitutes not merely a
signal processing pipeline but an alternative foundation for
machine intelligence. This is where the framework's ambitions
become most explicit and where the distance between engineering
results and architectural claims becomes most visible.
